\section{The one-variable expression and its witness maps}
\label{sec:reduction}

We combine the lattice gap with the rational density construction. Define
\begin{equation}
 \begin{gathered}
 h_j(t)=t\sin(t^{2j-1}),\qquad v_j(t)=h_j(t)^2,\\
 A_P(t)=F_P(v_1(t),\ldots,v_n(t)),\\
 E_P(t)=\abs{A_P(t)}-A_P(t).
 \end{gathered}\label{eq:compiled-expression}
\end{equation}
The map sending a polynomial code to this expression code is primitive
recursive: compute $C,D,n$, form the finite syntax trees in
\eqref{eq:majorant}--\eqref{eq:gap-function}, and substitute the indicated
one-variable expressions. Subtraction is multiplication by $-1$ followed by
addition. Thus $E_P$ uses only rational constants, $\pi$, the variable,
addition, multiplication, sine, and one outer occurrence of absolute value.
The inner expression $A_P$ is analytic. Exponential and $\log2$, although
allowed in the requested class, are not required by this reduction.

\subsection{A polynomial root produces a rational nonzero value}

Suppose $m\in\N^n$ satisfies $P(m)=0$. Compute $M,L_0,\rho$ as in
\eqref{eq:root-radius}, and put
\begin{equation}
 \eta=\frac{\rho}{8(M+1)}.\label{eq:square-tolerance}
\end{equation}
We first construct rational targets for the dense curve. For each coordinate
find $a_i\in[0,M+1]\cap\Q$ such that
\begin{equation}
 \abs{a_i^2-m_i}\le\rho/4.\label{eq:square-target}
\end{equation}
Here is an explicit construction that does not require a real square-root
operation. Bisect the rational interval $[0,M+1]$, maintaining endpoints
$l,u$ with $l^2\le m_i\le u^2$. The branch at each midpoint is determined
by an exact rational comparison. Stop after a prescribed number $N$ of
steps satisfying
\[
 2^N\ge\frac{8(M+1)^2}{\rho}.
\]
The resulting width is at most $\rho/[8(M+1)]$. Taking $a_i=l$ gives
\[
 0\le m_i-a_i^2\le u^2-l^2
 \le2(M+1)(u-l)\le\rho/4,
\]
as required. An integer ceiling of the binary logarithm of the displayed
bound supplies $N$ by a primitive recursive calculation; a larger elementary
bound is equally acceptable.

Apply Theorem~\ref{thm:density} to these rational targets, the integer bound
$A=M+1$, and tolerance $\eta$. It produces a rational $q$ with
$\abs{h_i(q)-a_i}<\eta$. Since $\eta\le1$ and $a_i\le M+1$,
\[
 \abs{h_i(q)^2-a_i^2}
 \le\eta(2\abs{a_i}+\eta)
 \le3(M+1)\eta=3\rho/8.
\]
Combining this with \eqref{eq:square-target} proves
\[
 \abs{v_i(q)-m_i}<5\rho/8<\rho.
\]
Each $v_i(q)$ is nonnegative, so Lemma~\ref{lem:root-neighborhood} applies:
\begin{equation}
 A_P(q)\le-3/4,
 \qquad
 E_P(q)\ge3/2.\label{eq:robust-nonzero}
\end{equation}

Recall that $\App(e,q,k)$ is rational and approximates the value of the
expression coded by $e$ at $q$ with error at most $2^{-k}$, and
\[
 \NZ(e,q,k)\;:\!\!\iff\;
 \abs{\App(e,q,k)}>2^{-k}.
\]
With $e$ the code of $E_P$, equation~\eqref{eq:robust-nonzero} gives
\[
 \App(e,q,3)\ge3/2-1/8=11/8>1/8.
\]
Thus $(q,3)$ is a strict nonzero-evaluation certificate. The margin is
deliberately generous. The witness construction never searches until an
unknown sign happens to become visible.

\subsection{A nonzero certificate produces a polynomial root}

Conversely, suppose $\NZ(e,q,k)$ holds for $e=E_P$. Put $d=2^{-k}$ and
$b=\App(e,q,k)$. The defining expression satisfies $E_P(q)\ge0$, so its
approximation satisfies $b\ge-d$. The rational comparison
$\abs b>d$ therefore implies $b>d$. Hence
\[
 \gamma=b-d>0,
 \qquad E_P(q)\ge\gamma.
\]

To obtain a negative value of $A_P(q)$ constructively, approximate it by a
rational $c$ to error at most $\tau=\gamma/8$. The required dyadic
precision is computed from the positive rational $\gamma$. For example,
the deliberately coarse integer choice
$r=\lceil8/\gamma\rceil+1$ ensures $2^{-r}\le\gamma/8$, since
$2^r\ge r+1$. Thus one may simply use $c=\App(A_P,q,r)$. If
$c\ge-\tau$, then $A_P(q)\ge-2\tau=-\gamma/4$. For any real $a\ge-r$
with $r\ge0$, the absolute-value triangle inequality gives
\[
 \abs a\le\abs{a+r}+r=a+2r,
 \quad\hbox{and therefore}\quad \abs a-a\le2r.
\]
Applying this with $r=\gamma/4$ would yield
$E_P(q)\le\gamma/2$, contradicting $E_P(q)\ge\gamma>0$.
Consequently the decidable rational comparison must give $c<-\tau$, and
then
\[
 A_P(q)\le c+\tau<0.
\]
This reasoning does not split on an undecidable sign of a real quantity; the
only branch is a rational comparison.

Now evaluate each $v_i(q)=h_i(q)^2$ to error at most $1/8$, and round the
result as in Lemma~\ref{lem:rounding}. The tuple $z$ obtained belongs to
$\N^n$. Since
$F_P(v_1(q),\ldots,v_n(q))=A_P(q)<0<1$, Theorem~\ref{thm:lattice-gap}
gives $P(z)=0$. This is a bounded computation from the certificate.
Indeed, the candidate tuple $z$ can be computed from $P$ and $q$ alone;
the supplied nonzero certificate is what proves that the candidate is a
root. No unbounded Diophantine search occurs.

For clarity, even the approximation of a squared coordinate has a direct
precision bound. Let $Q=1+\lceil\abs q\rceil$, so
$\abs{h_i(q)}\le Q$. Compute a rational $s_i$ approximating $h_i(q)$
to $\delta=1/[32(Q+1)]$. Then
\[
 \abs{s_i^2-h_i(q)^2}
 \le\delta(2Q+\delta)
 \le\frac{2Q+1}{32(Q+1)}<\frac18.
\]
Thus $s_i^2$ is an admissible rational input to the rounding rule. The
needed approximation to $h_i(q)=q\sin(q^{2i-1})$ is obtained by computing
the rational power exactly and approximating its sine to
$\delta/(Q+1)$. This spells out the reverse computation using only the
basic sine approximation primitive and rational arithmetic.

\subsection{The exact arithmetical reduction}

\begin{theorem}[Two finite witness transformations]\label{thm:witness-maps}
There are primitive recursive maps $\operatorname{forth}$ and
$\operatorname{back}$ with the following properties. For a well-formed
polynomial code $P$:
\begin{enumerate}
\item if $m\in\N^n$ and $P(m)=0$, then
      $\operatorname{forth}(P,m)=(q,3)$ satisfies $\NZ(E_P,q,3)$;
\item if $\NZ(E_P,q,k)$, then
      $z=\operatorname{back}(P,q,k)$ lies in $\N^n$ and satisfies
      $P(z)=0$.
\end{enumerate}
Their correctness is expressible, and provable, in $\ipra$ using coded
integer and rational operations and the finite approximation bounds above.
\end{theorem}

\begin{proof}
The preceding two constructions define the maps. On malformed codes one
may assign a fixed default output, making them total primitive recursive
functions. The first uses the explicit radius, bounded rational bisection,
and the density algorithm. The second uses fixed-precision approximations
and rational floor. All precision requirements and loop bounds are computed
before their associated finite loops.

For the arithmetical correctness assertions, the predicates $P(m)=0$ and
$\NZ(E_P,q,k)$ are primitive recursive relations on codes. Each displayed
analytic inequality is justified at the required precision by a finite
rational error bound. The sine and $\pi$ approximation identities can be
unfolded into those finite bounds, as described in the arithmetic semantics
and the appendix. All uses of induction concern bounded calculation records
or explicit primitive recursive recurrences. The integer conclusion in the
reverse map has a gap: $A_P(q)<0$ gives $H_P(v(q))<1/4$, and hence
$P(z)^2<1/2$. Thus the final integer comparison is separated from one by
a fixed rational margin. The real notation introduces no separate axiom
of real completeness, choice, or sign decidability.
\end{proof}

In the notation for solvability and identity used earlier, the two maps give
\begin{equation}
 \ipra\ \vdash\quad
 \Sol(P)\ \longleftrightarrow\ 
 \exists q\,\exists k\;\NZ(E_P,q,k).\label{eq:positive-reduction}
\end{equation}
Rational values $q$ here mean their natural-number codes. Since
$\Id(e)$ is defined as $\forall q\,\forall k\,\neg\NZ(e,q,k)$, pure
intuitionistic logic yields
\begin{equation}
 \boxed{\ \ipra\ \vdash\quad
          \Id(E_P)\ \longleftrightarrow\ \neg\Sol(P).\ }
 \label{eq:negative-reduction}
\end{equation}
For completeness, the two directions are immediate without a quantifier
negation principle. If $\Id(E_P)$ and a root $m$ are given, the forward
map produces a forbidden nonzero certificate. If $\neg\Sol(P)$ and a
nonzero certificate are given, the reverse map produces a forbidden root.
These are precisely the two implications in
\eqref{eq:negative-reduction}.

We have not inferred a positive nonzero certificate from the mere negation
of $\Id(E_P)$. In intuitionistic logic, that inference would require an
additional principle. The positive equivalence
\eqref{eq:positive-reduction} concerns \emph{witnesses}; the negative
equivalence \eqref{eq:negative-reduction} concerns their absence. Keeping
these assertions separate is essential to the final diagonal argument.

Externally, the connection with the usual real-valued identity is equally
direct. If $P$ has no natural-number root, every tuple $v(t)$ is
nonnegative, so the lattice gap gives $A_P(t)\ge1$ and $E_P(t)=0$ for
every real $t$. If $P$ has a root, the construction supplies a rational
$q$ with $E_P(q)\ge3/2$. Thus the conventional semantic identity and the
arithmetical predicate have the intended meaning on every expression used
by the reduction.
